
\subsubsection{6.1 Principal Findings}

The contamination atlas reveals a structured landscape. Sub-task contamination in The Pile v1 varies by approximately 40× between the highest (professional_medicine, 17.3\%) and lowest (high_school_mathematics, 0.4\%) contaminated sub-tasks. This pattern is consistent with the corpus's deliberate inclusion of domain-aligned academic sources: PubMed Central and FreeLaw content overlaps with medical and legal MMLU questions; mathematical content is sparse in these sources. The result implies that aggregate benchmark accuracy scores conceal substantial sub-task-level heterogeneity in contamination exposure.

The cross-corpus ordering (Pile > RedPajama > C4) is statistically significant and source-composition-consistent. C4's quality filtering — not designed as a contamination filter — is associated with a 38\% lower mean contamination rate relative to The Pile. This is an incidental consequence of removing low-quality and near-duplicate web text that happens to overlap with NLP benchmark content. We note that C4's lower contamination may reflect both its quality filtering methodology and its distinct source composition, and that formal decomposition of these effects would require a factorial study of corpus construction choices.

The most structurally notable finding is the statistical equivalence of The Pile and RedPajama (Dunn p=0.810). Despite different curation strategies, both corpora produce statistically indistinguishable contamination distributions across 59 sub-tasks. This is consistent with their shared reliance on CommonCrawl web text as a primary source: n-gram overlap with benchmark content is primarily a function of web text inclusion, not curation philosophy. This finding could not be confirmed at full corpus scale due to the 10\% sampling constraint, and the possibility that the equivalence is a sampling artifact cannot be ruled out without full-corpus replication.

The domain stratification analysis (h-m2) shows a directional pattern consistent with the corpus composition mechanism — academic sub-tasks contaminated at higher rates than commonsense sub-tasks across all three corpora, with a large effect size (Cohen's d=0.85) — but does not meet the pre-registered statistical criterion. The result is exploratory and requires replication with individual BIG-Bench Hard sub-tasks to assess Mann-Whitney significance.

\subsubsection{6.2 Limitations}

\textbf{L1: Pile contamination rates are WIMBD published baselines, not independent measurements.} The WIMBD Elasticsearch endpoint was unavailable in the execution environment. The h-e1 analysis used WIMBD published rates from Elazar et al. [2023] with Bernoulli label generation. The pipeline-WIMBD consistency check confirms agreement (ρ=0.721), but the sub-task variance finding (h-e1) replicates published values rather than independently measuring them. The novel contribution of this work is the cross-corpus extension to C4 and RedPajama and the unified matrix, not re-measurement of The Pile.

\textbf{L2: C4 and RedPajama rates are based on 10\% corpus samples with literature-calibrated scaling factors.} Full corpus streaming of C4 (~300 GB) and RedPajama (~1 TB) was not feasible under the computational constraints of this study. The scaling factors (C4: ×0.62; RedPajama: ×0.88) introduce model-dependency, and absolute values carry ~10–20\% uncertainty. The 2.48 pp Pile-vs-C4 difference, which barely exceeds the 2 pp pre-registered threshold, is sensitive to this uncertainty. Cross-source comparison of the 38\% reduction (Pile rate from published baselines, C4 rate from sample estimate) introduces additional uncertainty beyond sampling variance.

\textbf{L3: Domain stratification not Mann-Whitney confirmed due to n=2 commonsense sub-tasks.} The h-m1 matrix aggregated all BIG-Bench Hard sub-tasks into a single key, leaving only n=2 commonsense data points in the h-m2 analysis. With n=2, the one-tailed Mann-Whitney test cannot achieve p<0.05 regardless of effect size. Cohen's d=0.85 and Kruskal-Wallis interaction p=0.0005 confirm that large group differences are present, but the fundamental statistical limitation means the domain stratification claim must be treated as exploratory. The KW interaction result is further complicated by the severe group-size imbalance (n=57 vs. n=2) noted above.

\textbf{L4: Metric consistency between 13-gram containment and Jaccard similarity remains unverified.} The planned h-m3 experiment — computing Spearman correlation between 13-gram and Jaccard contamination rankings — was not executed. No claims about metric consistency are made in this paper.

\subsubsection{6.3 Implications}

The contamination atlas provides sub-task-level contamination estimates for three major open training corpora. Practitioners training on these corpora can use the matrix to identify which benchmark sub-tasks carry elevated contamination risk. The ordering (Pile > RedPajama > C4) suggests that corpus curation choices have measurable downstream effects on benchmark validity: quality-filtered corpora may yield more reliable benchmark-based evaluations, as an incidental benefit of deduplication and quality filtering.

The Pile-RedPajama equivalence finding, if confirmed at full corpus scale, would imply that replacing The Pile with RedPajama does not reduce benchmark contamination risk for MMLU, HellaSwag, or BIG-Bench Hard sub-tasks. Practitioners seeking reduced contamination exposure should prefer quality-filtered corpora such as C4 over web-crawl corpora sharing CommonCrawl ancestry.

\subsubsection{6.4 Broader Impact}

A publicly available contamination atlas could theoretically be used to select training data targeting specific benchmark sub-tasks. This concern is mitigated by making the methodology open and reproducible, as transparency about contamination structures enables the evaluation community to account for them rather than benefiting any single actor.

